\chapter{Discussion}
The results indicate that LDCT MOS prediction is primarily limited by task calibration rather than by the availability of generic medical visual features. Although the base MedGemma model provides a strong medical VLM prior, direct prompting does not produce reliable scalar MOS estimates. In contrast, SFT exposes the model to the required image-conditioned score format and aligns the generated response with the radiologist MOS scale. This explains why supervised adaptation improves both absolute score accuracy and rank consistency while preserving complete prediction coverage.

The SFT ablation suggests that this calibration is not simply limited by adapter capacity. Once the model learns the required response format and MOS scale, compact LoRA updates appear sufficient to capture the dominant quality cues in the evaluated split. The modest variation across SFT configurations is important for LDCT IQA, where training data are limited and over-parameterized adaptation can increase the risk of fitting score-specific artifacts rather than general image-quality cues.

GRPO behaves differently from SFT because the reward provides feedback on numerical closeness but does not by itself teach the response format or preserve the full supervised score distribution. This helps explain why reward optimization is most useful after supervised calibration, yet still remains below the strongest SFT models. The policy update can shift probability mass toward a small set of reward-favorable scores, reducing calibration and dynamic range. The banded predictions in Figure~\ref{fig:grpo-selected-scatter} are consistent with this interpretation. For MOS-based IQA, such compression is undesirable because clinically meaningful quality differences may be subtle and ordinal.

The GRPO sweep also shows that generation settings are part of the optimization problem, not only decoding details. Completion length, sampling temperature, candidate-group size, and KL regularization affect whether the model can keep producing parseable scalar answers while being updated by the reward. For scalar MOS prediction, controlling output validity and score calibration is therefore at least as important as optimizing the reward itself.

These findings support using SFT as the primary adaptation stage for LDCT MOS prediction and treating GRPO as an optional refinement that requires careful reward and sampling design. Future work should compare the generative formulation with direct regression heads that predict MOS as a scalar output, and should analyze performance by anatomy, dose level, and MOS range. Another promising direction is to extend the task beyond scalar MOS prediction toward reasoning-oriented image-quality assessment, where the model predicts a score together with clinically grounded explanations of noise, artifacts, contrast loss, and diagnostic acceptability.
